\begin{frame}{Lecture summary}
  \small
  \begin{itemize}\itemsep0.22em
    \item[\textcolor{red}{$\blacktriangleright$}] \textbf{Motivation:} video is the natural domain for dynamics and ``world'' hypotheses; we contrast \textbf{pixel generative} stacks with \textbf{latent predictive} models.
    \item[\textcolor{red}{$\blacktriangleright$}] \textbf{World Models:} Can be used to generate interactive videos, 3D scenes and latent world representations.
    \item[\textcolor{red}{$\blacktriangleright$}] \textbf{The Simulation Gap:} Despite stunning visuals (e.g., pouring liquid), current models lack symbolic or true physics-based understanding (e.g., accurate fluid dynamics).
    \item[\textcolor{red}{$\blacktriangleright$}] \textbf{Takeaway:} The two families answer different questions. Ask a pixel model for a convincing clip; ask a latent model for a state you can plan against. Which one you need is set by the task, not by sample quality.
    \item[\textcolor{red}{$\blacktriangleright$}] \textbf{Looking ahead:} The open problems are duration and controllability: rollouts stay coherent for minutes rather than hours, and action sets remain closed while the space of requested world changes is open.
  \end{itemize}
\end{frame}
